\documentclass[aspectratio=169,10pt,spanish]{beamer}

\input{../../../_preambulo_beamer.tex}
\input{../../../_comandos_beamer.tex}

\selectlanguage{spanish}

\title{MA1001B - Modelación Estadística para la Toma de Decisiones}
\subtitle{Lección 11: Teorema de Bayes y actualización de probabilidad}
\author{}
\institute{Tecnol\'ogico de Monterrey}
\date{}

\begin{document}

\begin{frame}[plain]
  \vspace{1.2cm}
  {\Large\bfseries MA1001B - Modelación Estadística para la Toma de Decisiones\par}
  \vspace{0.7cm}
  {\large Lección 11: Teorema de Bayes y actualización de probabilidad\par}
  \vspace{0.7cm}
  {\color{TecRojo}\rule{\textwidth}{1.2pt}}
  \vspace{0.8cm}

  Facultad MA1001B\\[0.3cm]
  Periodo\\[0.3cm]
  Tecnol\'ogico de Monterrey
\end{frame}

\begin{frame}{De una condicional inversa a una actualización}
  \begin{alertblock}{Pregunta guía}
    ¿Cómo debe actualizarse una tasa priori de defectuosos después de una
    marca de inspección, y qué ocurre si esa priori es mucho menor que
    $0.10$?
  \end{alertblock}

  \vspace{0.6em}
  Al final de esta lección, usted puede:
  \begin{itemize}
    \item separar la priori $P(D)$ de las verosimilitudes $P(F\mid D)$ y $P(F\mid C)$;
    \item calcular $P(D\mid F)$ con el teorema de Bayes;
    \item hacer coincidir la posterior con la razón de caminos $8/17$;
    \item explicar por qué $P(F\mid D)$ no es una posterior cuando se ignora la
      tasa base.
  \end{itemize}

  \vfill
  \scriptsize Todos los lotes y marcas de inspección usados hoy son sintéticos. No se usan datos reales de empresa.
\end{frame}

\begin{frame}{Priori, verosimilitud, evidencia}
  \small
  Continúe el árbol de la Lección 10 con $P(D)=0.10$, $P(F\mid D)=0.80$ y
  $P(F\mid C)=0.10$.
  \[
    P(D\mid F)=\frac{P(F\mid D)\,P(D)}{P(F\mid D)\,P(D)+P(F\mid C)\,P(C)}.
  \]

  \vspace{0.4em}
  \begin{center}
    \begin{tabular}{lcc}
      \toprule
      \textbf{Pieza} & \textbf{Símbolo} & \textbf{Rol} \\
      \midrule
      Priori & $P(D)=0.10$ & Tasa antes de la marca \\
      Verosimilitud & $P(F\mid D)=0.80$ & Detección dado defectuoso \\
      Tasa de marca falsa & $P(F\mid C)=0.10$ & Marca dado limpio \\
      Evidencia & $P(F)$ & Probabilidad total de una marca \\
      \bottomrule
    \end{tabular}
  \end{center}
\end{frame}

\begin{frame}[fragile]{Paso 1: Priori y verosimilitudes}
  \begin{columns}[T]
    \begin{column}{0.42\textwidth}
      \small
      \begin{lstlisting}
prior = 0.10
p_f_given_d = 0.80
p_f_given_c = 0.10
      \end{lstlisting}

      \begin{block}{Salida verificada}
        $P(D)=0.100000$\\
        $P(C)=0.900000$\\
        $P(F\mid D)=0.800000$\\
        $P(F\mid C)=0.100000$
      \end{block}
    \end{column}
    \begin{column}{0.58\textwidth}
      \centering
      \includegraphics[width=\textwidth]{../figuras/bayes_01_priori_verosimilitud.png}
      \vspace{0.2em}
      \scriptsize Priori versus verosimilitudes de marca del Paso 1
    \end{column}
  \end{columns}
\end{frame}

\begin{frame}{Bayes combina los dos caminos marcados}
  \small
  Numerador: el camino defectuoso y marcado
  \[
    P(D\cap F)=0.80\times 0.10=0.08.
  \]
  Denominador: ley de la probabilidad total
  \[
    P(F)=0.08+0.90\times 0.10=0.17.
  \]
  Posterior:
  \[
    P(D\mid F)=\frac{0.08}{0.17}=\frac{8}{17}=0.470588.
  \]
\end{frame}

\begin{frame}[fragile]{Paso 2: Posterior a partir del teorema de Bayes}
  \begin{columns}[T]
    \begin{column}{0.40\textwidth}
      \small
      \begin{lstlisting}
num = p_f_d * p_d
den = num + p_f_c * p_c
post = num / den
      \end{lstlisting}

      \begin{block}{Salida verificada}
        Numerador $=0.080000$\\
        $P(F)=0.170000$\\
        $P(D\mid F)=0.470588$\\
        Coincide con $8/17$
      \end{block}
    \end{column}
    \begin{column}{0.60\textwidth}
      \centering
      \includegraphics[width=\textwidth]{../figuras/bayes_02_posterior.png}
      \vspace{0.2em}
      \scriptsize Numerador, $P(F)$ y posterior del Paso 2
    \end{column}
  \end{columns}
\end{frame}

\begin{frame}{Paso 3: Una priori rara colapsa la posterior}
  \begin{columns}[T]
    \begin{column}{0.42\textwidth}
      \small
      Mantenga $P(F\mid D)=0.80$ y $P(F\mid C)=0.10$. Cambie solo $P(D)$ a
      $0.02$.

      \vspace{0.4em}
      \[
        P(F)=0.114000,
      \]
      \[
        P(D\mid F)=0.140351.
      \]

      \vspace{0.4em}
      \textbf{Límite:} tratar $P(F\mid D)=0.80$ como posterior ignora la tasa
      base. La Lección 12 pregunta si los conteos mismos vinieron de un marco
      sesgado.
    \end{column}
    \begin{column}{0.58\textwidth}
      \centering
      \includegraphics[width=\textwidth]{../figuras/bayes_03_sensibilidad_tasa_base.png}
      \vspace{0.2em}
      \scriptsize Posterior versus priori del Paso 3
    \end{column}
  \end{columns}
\end{frame}

\begin{frame}{Verificación de conceptos}
  \begin{enumerate}
    \item ¿Cuál número es la priori y cuál es una verosimilitud: $0.10$ o
      $0.80$?
    \item Calcule $P(F)$ cuando $P(D)=0.10$.
    \item ¿Por qué $P(D\mid F)=8/17$ coincide aquí con el teorema de Bayes?
    \item Si $P(D)=0.02$, ¿por qué $P(D\mid F)$ queda muy por debajo de
      $0.80$?
  \end{enumerate}

  \vspace{0.6em}
  \begin{block}{Conexión con la Lección 12}
    La sesión puede continuar directamente con sesgo de muestreo, huecos de
    cobertura y error no muestral en un registro sintético.
  \end{block}

  \vfill
  \scriptsize Fuente: OpenStax, \textit{Principles of Data Science},
  Sección 3.4; \textit{Introductory Business Statistics 2e}, Sección 3.4.
  CC BY-NC-SA.
\end{frame}

\appendix



\end{document}
